\documentclass[11pt]{article}
\usepackage[margin=1in]{geometry}
\usepackage{amsmath,amssymb,amsthm,mathtools}
\usepackage[hidelinks]{hyperref}
\setlength{\emergencystretch}{3em}
\allowdisplaybreaks[2]
\newtheorem{theorem}{Theorem}[section]
\newtheorem{lemma}[theorem]{Lemma}
\theoremstyle{remark}\newtheorem{remark}[theorem]{Remark}
\newcommand{\dd}{\,d}
\newcommand{\one}{\mathbf1}
\newcommand{\eps}{\varepsilon}
\newcommand{\Kfun}{\mathcal K}
\title{Secondary Q2 entry profile at strong half mass}
\author{Analytical increment for review}
\date{Version 1; October 21, 2026}
\begin{document}
\maketitle
\begin{abstract}
For the exact original initialized flow, the whole fixed Q2 group
$c_0<Kq^\gamma$ has a bounded-displacement entry profile about
the actual strong reference at the first strong half-mass time.
Its raw radial mass is of order $A_q^{-1}q^3$ and its secondary
tail has index $(1-3\lambda/2)/d$, with displacement scale
$A_q^d q^{10d-2s}$. The angular comparison has fixed two-sided
constants; no universal asymptotically exact angular multiplier
is asserted. Same-field subtraction makes the chart error
homogeneous in displacement, including for the smallest labels.
Only the entry profile is supplied. The later distance envelope
and capped gain remain separate S2 inputs.
\end{abstract}

\section{Setting and exact original-flow result}
All new labels have prefix \texttt{an02sepv1:}.
Use the model and initialization of MOMC and Q2EC:
\[
 P=\tfrac12N(e_1,q^2I_2)+\tfrac12N(\rho e_2,q^2I_2),\qquad
 \rho\in[13/20,7/10],\qquad
 U(0,\beta)=W(0,\beta)=q^5(\cos\beta,\sin\beta),
 \quad\dd\lambda_0=\dd\beta/(2\pi).
\]
Balance gives $U=\sqrt m\,a_\theta$, $W=\sqrt m\,a_\psi$.
Put $\lambda=\rho^2$, $\omega=(1-\lambda)/2$, $\ell=\log(1/q)$,
\[
 \Kfun(z)=\phi(z)+z\Phi(z),\quad a=\rho\Kfun(\rho a),\quad
 \alpha=\lambda\Phi(\rho a)/2,\quad d=1/2-\alpha,\quad s=d/\omega,
 \qquad \gamma=1-5\lambda+2\lambda/(1-\lambda).
\]
PROFILE's exact target-only root is $a_q=a+O(q^2)$.
Hats always denote the exact target-only characteristic with
the same initial label and at the same physical time.
Let $M$ be the mass of the fixed initial right half-circle.
Use MOMC's proved clock, without changing its normalization:
\begin{align}
 t_m&=(2/\omega+1/2)\ell,\quad
 u_m=10-2/\omega-1/2,\quad v_m=u_m-1,\nonumber\\
 M_0&=M(t_m)=A_q q^{u_m},\quad c_A\le A_q\le C_A,\nonumber\\
 t_0=t_h&=10\ell-\log A_q+O(q^2\ell^2+q^{u_m}),\quad M(t_h)=1/2.
 \label{an02sepv1:clock}
\end{align}
Here $u_m\ge169/102$ and $v_m\ge67/102$.
The actual reference is the characteristic launched at
$\beta_c=qa_q$. Write
\[
 C(t)=(\theta_c(t)/q,\psi_c(t)/q),\qquad
 V_\beta(t)=((\theta_\beta-\theta_c)/q,
                    (\psi_\beta-\psi_c)/q),\qquad
 \xi_\beta=\|V_\beta(t_h)\|_1.
\]
In particular $C$ is never reset to $(a,a)$ or $(a_q,a_q)$.
MOMC gives $|C_i|\le R=79/51$ up to $t_h$.

Fix Q2EC's $C_0$, and its already chosen $K=K(C_0)$ and $Z$.
For $\beta\in[\pi/2,\pi]$ put $c_0=-\cos\beta$, $s_0=\sin\beta$.
The fixed group and its raw entry measure are
\[
 \Sigma=\{\beta\in[\pi/2,\pi]:c_0<Kq^\gamma\},\qquad
 \eta(B)=\int_\Sigma m(t_h,\beta)\one_{\{\xi_\beta\in B\}}\dd\lambda_0.
\]
There is no normalization by a target-only strong mass and no
moving membership indicator in this definition. Define
\begin{equation}
 b=2d/\lambda,\quad a_2=2/\lambda-2,\quad
 p_2=(a_2-1)/b=\frac{1-3\lambda/2}{d},\quad
 h_A=A_q^d q^{10d-2s},\quad
 D_h=\sqrt{2(1-M_0)},\quad B_q=A_q^{-1}D_h^{-2}.
 \label{an02sepv1:scales}
\end{equation}
All $\asymp$ constants below are uniform in $q$ and $\rho$;
they can depend on the fixed $K,Z,c_A,C_A$ and proved source
constants. In particular they are independent of LCRC's guards.

\begin{theorem}[Recruited entry with fixed constants]
\label{an02sepv1:main}
For sufficiently small $q$, every label in $\Sigma$ satisfies
$V_{\beta,1}(t_h)>0$, $V_{\beta,2}(t_h)>0$, and
\begin{align}
 \xi_\beta&\le Z_\Sigma:=Z+2R,\qquad
       \eta([0,Z_\Sigma])\le C_{\rm ent} A_q^{-1}q^3,\label{an02sepv1:cap}\\
 V_{\beta,i}(t_h)&\asymp h_A
       \left[\frac{c_0+q}{q(s_0+q)}\right]^b\quad(i=1,2),
       \label{an02sepv1:labelshape}\\
 m(t_h,\beta)&\asymp B_q q^{2/\lambda}
       (c_0+q)^{2-2/\lambda}(s_0+q)^{2/\lambda}.
       \label{an02sepv1:labelweight}
\end{align}
One can retain the more precise radial comparison
\begin{equation}
 \frac{m(t_h,\beta)}{\widehat m(t_h,\beta)}
   =D_h^{-2}\exp\{O(\mathcal R_q)\},\qquad
 \mathcal R_q=q^{1/1000}+q^{4/51}+q^{v_m}+q\ell^2+q^2\ell^2=o(1).
 \label{an02sepv1:radialfactor}
\end{equation}
The constants in this error are allowed to depend on the fixed
receiver chart used in the proof.
The total mass is in fact $\asymp B_q q^3$. For every $w>0$,
\begin{equation}
 \eta\{\xi>w\}\le C B_q q^3
             \min\{1,(h_A/w)^{p_2}\}.
 \label{an02sepv1:tailupper}
\end{equation}
There are fixed $C_1,z_0>0$ such that
\begin{equation}
 \eta\{\xi>w\}\asymp B_q q^3(h_A/w)^{p_2}
       \quad(C_1h_A\le w\le z_0).
 \label{an02sepv1:taillower}
\end{equation}
This is Q2E's interior tail range; no lower tail is asserted
arbitrarily close to the last-label cutoff.
Since $A_q,D_h$ are bounded above and below, these estimates
supply LCRC's entry-profile contract with $h_1=q^{10d-2s}$.
\end{theorem}

One may take $C_{\rm ent}=4e^7C_Q$, where $C_Q$ is the
target energy constant specified in Q2EP Theorem 2.1.
The cap and its upper mass estimate in \eqref{an02sepv1:cap}
are immediate from Q2EC and Q2EP, respectively: subtract the
reference, whose coordinate sum is at most $2R$, and restrict
the whole-Q2 mass estimate to $\Sigma$. We do not reprove those
results. The work below proves the label laws and their tail.

\section{Two domain checks for the source comparisons}
\begin{lemma}[Early relative comparison on the required label arc]
\label{an02sepv1:early}
At $t_m$, for every $\beta\in[\pi/2,\pi]$,
\[
 V_{\beta,i}(t_m)=
   (\widehat\theta(t_m,\beta)/q-a_q)e^{O(q^{v_m})}\quad(i=1,2),
 \qquad m(t_m,\beta)/\widehat m(t_m,\beta)=e^{O(q^{u_m})}.
\]
This extends OWPT's early tangent argument to the arc
$[qa_q,\pi]$; it does not extend a target profile beyond its domain.
\end{lemma}
\begin{proof}
Here is the domain check, since OWPT states its pairwise lemma
on the right half-circle. P's early angle/radial comparison is
global in the initial label. With
$\overline S=q^{10}e^{(1+2q^2)t}$ it gives angular error
$O(\overline S)$ and radial logarithmic error $O(\overline S)$.
Every estimate in OWPT's receiver-Jacobian proof is global in
receiver angle: $0\preceq A_\theta\preceq a_*I$,
$A_\theta'a_\theta=0$, $\|A_\theta'\|\le32/q$, and
$\|DF^f\|\le C\overline S/q$ by PROFILE Proposition 4.1.
Consequently the same calculation gives, on the entire arc,
\[
 DF=\begin{pmatrix}-g&h\\h&-g\end{pmatrix}+E_\beta,
 \quad h\ge0,\quad \|E_\beta\|_\infty\le C\overline S/q.
\]
The positive exact target tangent $\widehat p$ solves
$\widehat p'=(h-g)\widehat p$. Thus $w=p/\widehat p$ solves
\[
 w'=h\begin{pmatrix}-1&1\\1&-1\end{pmatrix}w+E_\beta w,
 \qquad w(0)=\one.
\]
The first propagator is nonnegative, fixes $\one$, and has
$\ell^\infty$ norm one. Its variation-of-constants estimate
is $\|w-\one\|_\infty\le e^{\int\|E\|}-1$.
Now $\int_0^{t_m}\overline S/q=O(q^{v_m})=o(1)$.
Each tangent component is therefore in $e^{\pm Cq^{v_m}}$.
Integrating in the initial label from $qa_q$ to $\beta$, and
using the stationary target reference, proves the assertion.
This divides by a positive tangent, never by a small absolute
angle error. The radial assertion is P's early estimate.
\end{proof}

\begin{lemma}[Target profile on any fixed entry chart]
\label{an02sepv1:target}
At $t_h$, every $\beta\in\Sigma$ has positive
$\widehat h=\widehat\theta/q-a_q$, and
\begin{align*}
 \widehat h(t_h,\beta)&\asymp h_A
          [(c_0+q)/(q(s_0+q))]^b,\\
 \widehat m(t_h,\beta)&\asymp A_q^{-1}q^{2/\lambda}
          (c_0+q)^{2-2/\lambda}(s_0+q)^{2/\lambda}.
\end{align*}
Moreover every such target label is below $\pi/4$ at $t_m$,
for small $q$.
\end{lemma}
\begin{proof}
Q2EC's target-contract equation gives
$0<\widehat\theta(t_h,\beta)/q\le Z_T(K)$.
Choose a fixed chart threshold $L_T>Z_T(K)+a+1$, uniformly
over the ratio box. Q2E's profile proof remains valid for
this enlarged fixed chart and a bounded shift of $10\ell$.
To verify this assertion rather than importing its displayed
$T=10\ell$ literally, PROFILE's expansion
$f_q=f_0+O_{C^2}(q^2)$ holds on every fixed interval, where
\[
 f_0(x)=\tfrac12[\rho\Kfun(\rho x)-x],\quad
 f_0'\le-\omega,\quad f_0'(a)=-d.
\]
From the first hit of $qL_T$ write
$\widehat h'=-b_q(\widehat h)\widehat h$.
On this fixed chart $b_q\ge\omega/2$ and
$|b_q(\widehat h)-d|\le C_{L_T}\widehat h+C_{L_T}q^2$.
Thus $\int\widehat h\le2(L_T-a_q)/\omega$ and
\begin{equation}
 \widehat h(t)\asymp e^{-d(t-T_{L_T})},\qquad
 \widehat m(t)\asymp q^2\widehat m(T_\times)e^{t-T_\times}
 \quad(T_{L_T}\le t\le t_h).
 \label{an02sepv1:targetchart}
\end{equation}
The arrival formula is still
$T_{L_T}=T_\times+(2/\omega)\ell+O_{L_T}(1)$:
PROFILE's outer integral may be stopped at any larger fixed
$L_T$, changing only its bounded endpoint term. Its axis-start
first-coordinate and gate-integral proof gives the radial
formula in \eqref{an02sepv1:targetchart}. These observations
check the chart-size and time domains of Q2E, not a new
unproved comparison assumption.

Inserting Q2E's endpoint-regularized formulas
\[
 T_\times=\frac2\lambda\log\frac{c_0+q}{q(s_0+q)}+O(1),\qquad
 \widehat m(T_\times)\asymp q^{10}(c_0+q)^2/q^2
\]
gives its two label laws with $q^{-2s}e^{-dt_h}$ and
$q^{10}e^{t_h}$ in place of $q^{10d-2s}$ and one.
Equation \eqref{an02sepv1:clock} gives precisely the stated
$A_q^d$ and $A_q^{-1}$ factors. The finite-noise exponent
replacement costs $O(q^2\ell)$ in logarithms. Scalar uniqueness
keeps $\widehat h>0$ for these initial labels.

Finally $s_0\asymp1$ on $\Sigma$, and
$c_0+q\le (K+1)q^\gamma$, since $0<\gamma<1$.
The same crossing formula gives
$T_\times\le(10-2/\omega)\ell+O_K(1)$.
PROFILE's axis-to-$\pi/4$ part of the arrival proof takes
$\omega^{-1}\ell+O(1)$. Hence all these labels are below
$\pi/4$ by $(10-1/\omega)\ell+O_K(1)$, earlier than $t_m$:
\[
 3/\omega-19/2\ge 2400/231-19/2=137/154>0.
\]
No new late weak-axis crossing estimate is needed for $\Sigma$.
\end{proof}

\section{Homogeneous chart comparison about the actual reference}
Let $f^M$ be the coherent leading receiver field with strong
donors at $(a,a)$ and mass $M$:
\begin{align}
 2f_x^M(x,y)&=-(1-M)x+\rho\phi(\rho x)
       -\rho M F(\rho x,\rho a)+\lambda y\Phi(\rho x),\nonumber\\
 2f_y^M(x,y)&=-(1-M)y-Ma+\rho\Kfun(\rho x),\nonumber\\
 F(t,r)&=\phi(\min(t,r))+r\Phi(\min(t,r)).
 \label{an02sepv1:coherent}
\end{align}
This is a comparison vector field; it is not substituted for
the actual reference. It vanishes at $(a,a)$.
Its two off-diagonal receiver derivatives are
$\lambda\Phi(\rho x)/2\ge0$. On the positive aligned ray,
\begin{equation}
 f^M(a+h,a+h)=g_M(h)\one,\qquad
 g_M(h)=f_0(a+h)+Mh/2\le-h/200\quad(h\ge0,M\le1/2).
 \label{an02sepv1:ray}
\end{equation}
Indeed $F(\rho(a+h),\rho a)=\Kfun(\rho a)$ and
$\omega-M/2\ge51/200-1/4=1/200$.
The Jacobian at the center is
\[
 J_{\rm sh}(M)=\begin{pmatrix}-(1-M)/2&\alpha\\
                      \alpha&-(1-M)/2\end{pmatrix},\qquad
 r_0(M)=-d+M/2.
\]
In particular the $y$-diagonal is $-(1-M)/2$.

\begin{lemma}[Relative, rather than additive, chart error]
\label{an02sepv1:c1}
On any fixed receiver chart containing $C(t)$ and $C(t)+V$,
\begin{equation}
 V'=f^{M(t)}((a,a)+V)-f^{M(t)}(a,a)+E(t,V)V,
 \qquad \|E(t,V)\|_\infty\le\varepsilon_B(t),
 \label{an02sepv1:homogeneous}
\end{equation}
where $\sup\varepsilon_B=o(1)$ and
\begin{equation}
 \int_{t_m}^{t_h}\varepsilon_B\le C_B\mathcal E_q=o(1),\qquad
 \mathcal E_q=q^{v_m}+q^{s+d/2}+q^2(1+\ell)+q\ell^2.
 \label{an02sepv1:error}
\end{equation}
\end{lemma}
\begin{proof}
Let $e=\|C-(a,a)\|_\infty$, and let $\overline D$ be MOMC's
normalized first moment of donor displacement from $C$.
MOMC proves $\sup(e+\overline D)=o(1)$ and
$\int(e+\overline D)\le C\mathcal E_q$.
Its complementary ancestry mass obeys
$\mu\le Cq^2\ell^2M$ and $\int\mu/q\le Cq\ell^2$.
The finite-noise reduction is valid for any fixed finite receiver
radius, not only $79/51$; the donor chart remains $R=79/51$.
It bounds the value and first receiver derivative error by
$C_Bq^2+160\mu/q$.

Replacing the donor measure by $M\delta_{(a,a)}$ costs
$C_B M(e+\overline D)$ in $C^1$. Explicitly, the output donor
moment changes by that amount, and both $F(t,r)$ and
$F_t(t,r)=(r-t)_+\phi(t)$ are uniformly Lipschitz in $r$ on
the compact chart. This also handles equality of thresholds.
The coherent field is $C^{1,1}$ there.
Subtract the actual receiver velocity at $C$ from that at
$C+V$, and apply the mean-value formula to their difference.
Then shift both arguments of the coherent derivative by
$C-(a,a)$. This proves \eqref{an02sepv1:homogeneous} with
\[
 \varepsilon_B=C_B(q^2+\mu/q+e+\overline D).
\]
Its supremum and integral have the asserted bounds. Crucially
there is no additive $O(q^2+\mu/q+e)$ forcing in the equation
for $V$.
\end{proof}

\begin{lemma}[Positive displacement with bounded nonlinear distortion]
\label{an02sepv1:chartlemma}
Fix $H_*>0$ and $c_*>0$. Suppose at a time
$t_a\in[t_m,t_h]$ the actual displacement has
\[
 0<c_* H_a\le V_i(t_a)\le H_a\le H_*\quad(i=1,2).
\]
Choose a fixed chart larger than $R+H_*+1$.
Then both components stay positive to $t_h$, and, uniformly
for $t_a\le t\le t_h$,
\begin{equation}
 c_*'H_a e^{-d(t-t_a)+\frac12\int_{t_a}^tM}
 \le V_i(t)\le
 C_*'H_a e^{-d(t-t_a)+\frac12\int_{t_a}^tM}.
 \label{an02sepv1:charttransport}
\end{equation}
The positive constants depend only on the fixed chart,
$H_*,c_*$ and the ratio box. They do not require $H_a=o(1)$.
\end{lemma}
\begin{proof}
Stop at the first loss of positivity or exit from the fixed
chart. Put $H(t)=\max_i V_i(t)$. Cooperativity of
\emph{the coherent field}, not of the exact finite-noise field,
and \eqref{an02sepv1:ray} imply at a maximal component
\[
 D^+H\le-H/200+\varepsilon_B H\le-H/400.
\]
Thus $H\le H_a e^{-(t-t_a)/400}$ and
$\int_{t_a}^t H\le400H_*$. The chart cannot be exited.
Set
\[
 A(t)=\int_0^1 Df^{M(t)}((a,a)+rV(t))\dd r,\qquad
 P_0(t)=\exp[-d(t-t_a)+\tfrac12\int_{t_a}^t M].
\]
The matrix $A(t)$ is Metzler. Its row sums differ from $r_0(M)$
by at most $C_BH(t)$, by the local Lipschitz bound on $Df^M$.
The propagator $Q(t,s)$ for $A-r_0I$ is therefore positive and
\begin{equation}
 e^{-\int_s^t C_BH}\one\le Q(t,s)\one
       \le e^{\int_s^t C_BH}\one,
 \qquad\|Q(t,s)\|_\infty\le e^{\int_s^t C_BH}.
 \label{an02sepv1:positiveprop}
\end{equation}
These bounds follow by comparison with scalar multiples of
$\one$ in a cooperative linear system. They are analytic
row-sum bounds, not assumptions about exact-flow order.

Let $W=V/P_0$, $L_*=400C_BH_*$, and
$E_* =\int_{t_a}^{t_h}\varepsilon_B=o(1)$.
Variation of constants in $W'=(A-r_0I+E)W$ gives
\[
 \|W(t)\|_\infty\le H_a e^{L_*+E_*},\qquad
 \|W(t)-Q(t,t_a)V(t_a)\|_\infty
       \le H_a E_*e^{2L_*+E_*}.
\]
The unperturbed vector has each component at least
$c_*H_a e^{-L_*}$. For small $q$ the last error is less than
half this number. Hence $W_i\ge c_*H_a e^{-L_*}/2>0$,
uniformly up to the stopped endpoint. This excludes loss of
positivity and proves continuation. One may take
$c_*'=c_*e^{-L_*}/2$, $C_*'=e^{L_*+1}$.
All constants are fixed before choosing $q$. Large but finite
$e^{L_*}$ is allowed at the accuracy required here.
\end{proof}

\section{Matching every secondary label and its radial weight}
\begin{proof}[Proof of the label laws in Theorem \ref{an02sepv1:main}]
Choose a fixed $L$ at least as large as Q2EC's off-chart
threshold and $Z_T(K)+1$, and sufficiently large that
\[
 L\ge4(C_\Delta+R+a_{\max}+2),\qquad a_{\max}=28/51.
\]
Here $C_\Delta$ is Q2EC's fixed constant in
$|\psi-\theta|\le C_\Delta[q^{u_m}e^{-c_\Delta(t-t_m)}+qM(t)]$.
Increasing this auxiliary receiver threshold does not change
the already fixed group $\Sigma$ or Q2EC's $K$.
Use $H_*=8L$, $c_*=1/8$ and receiver radius
$B=R+8L+2$ in Lemma \ref{an02sepv1:chartlemma}.

If $\widehat\theta(t_m,\beta)\le2qL$, let $t_a=t_m$.
Lemma \ref{an02sepv1:early} gives positive initial components,
each equal to $\widehat h(t_m,\beta)e^{O(q^{v_m})}$.
Thus the chart lemma applies even when that displacement is
much smaller than the absolute reference offset.
The exact target is already in a fixed chart, where the
logarithmic decay calculation in Lemma \ref{an02sepv1:target}
gives
\[
 \widehat h(t,\beta)\asymp
       \widehat h(t_a,\beta)e^{-d(t-t_a)}.
\]
Its distortion constant is uniform in the initial displacement,
since its nonlinear integral is bounded by its initial size.

Otherwise the target and original input angles exceed $qL$
at $t_m$. By Lemma \ref{an02sepv1:target}, the target is already
below $\pi/4$ there. P's early comparison supplies positive
Cartesian matches with relative error $O(q^{v_m})$; in
particular the first-coordinate margin is fixed and the
second-coordinate margin is at least $qL$.
Apply Q2EC's two-sided off-chart transport with $t_s=t_m$
until the first time either input angle hits $qL$.
Such a $t_a\le t_h$ exists because the target endpoint is at
most $Z_T(K)q<Lq$.
The cotangent ratio in $[1/(2\sqrt2),2]$ implies that at this
first hit both input angles lie in $[qL,4qL]$ for small $q$.
The difference estimate and $|C_i|\le R$ put each component
of $V(t_a)$ between $L-C_\Delta-R-1$ and
$4L+C_\Delta+R+1$. These bounds imply the chosen
$H_*,c_*$ conditions, and also
\[
 V_i(t_a)\asymp\widehat h(t_a,\beta)\asymp L.
\]
The actual and target receivers are now on fixed charts;
apply the preceding chart comparisons.

In both cases we have proved the useful intermediate statement
\begin{equation}
 V_i(t_h)\asymp\widehat h(t_h,\beta)
      \exp\left(\frac12\int_{t_a}^{t_h}M\right).
 \label{an02sepv1:matchedshape}
\end{equation}
MOMC's mass clock gives
\begin{equation}
 \int_{t_m}^{t_h}M=\log[2(1-M_0)]+O(q^2\ell^2).
 \label{an02sepv1:massintegral}
\end{equation}
Thus the exponential in \eqref{an02sepv1:matchedshape} is
between $1$ and $\sqrt2\,e^{O(q^2\ell^2)}$.
All additional nonlinear factors in that equation are bounded
two-sided constants, as proved in the chart lemma.
Lemma \ref{an02sepv1:target} proves \eqref{an02sepv1:labelshape},
and positivity has been proved about the actual reference.

For completeness the radial comparison retains more accuracy.
In the early-chart case, P gives
$\log(m/\widehat m)(t_a)=O(q^{u_m})$.
In the late-chart case, Q2EC's first-coordinate estimate gives
\[
 \log\frac{U_1(t_a)}{\widehat U_1(t_a)}
       =-\tfrac12\int_{t_m}^{t_a}M+O(\mathcal R_q).
\]
At the hit, both angles are in a fixed scaled chart, so
$m=U_1^2/\cos^2\theta$ and its target version imply
\[
 \log\frac{m(t_a)}{\widehat m(t_a)}
       =-\int_{t_m}^{t_a}M+O(\mathcal R_q).
\]
After $t_a$, MOMC's radial reduction on the enlarged fixed chart
and its target version give
\[
 (\log m)'=1-M+O(C_Bq^2+3\mu),\qquad
 (\log\widehat m)'=1+O(C_Bq^2).
\]
The integrated error is $O(C_Bq^2\ell^2)$ because
$\int\mu\le Cq^2\ell^2\int M$.
Combining the two pieces and \eqref{an02sepv1:massintegral}
proves \eqref{an02sepv1:radialfactor}, uniformly in the label.
Inserting the target weight law proves
\eqref{an02sepv1:labelweight}. No post-$t_h$ gain law was used.
\end{proof}

\section{The tail, its cutoff, and the LCRC handoff}
\begin{proof}[Proof of the measure assertions]
Use the exact variable $y=1+c_0/q$ on $\Sigma$.
Its range is $1\le y<Y_q=1+Kq^{\gamma-1}$, and
$\dd\beta=q\dd y/s_0$.
Since $\gamma>0$, $s_0\asymp1$ uniformly on this group.
The two label laws yield fixed positive constants
$c_h,C_h,c_w,C_w$ such that
\begin{align}
 c_h h_A y^b&\le\xi\le C_h h_A y^b,\nonumber\\
 c_w B_q q^3y^{-a_2}\dd y
    &\le m(t_h,\beta)\dd\lambda_0
       \le C_w B_q q^3y^{-a_2}\dd y.
 \label{an02sepv1:density}
\end{align}
These are inequalities for the density in the initial variable;
monotonicity of the actual final displacement is not needed.
The endpoint regularizations $(c_0+q)$ and $(s_0+q)$ were
retained until this change of variables on the shrinking group.

Here $a_2>1$ uniformly. Integration over $[1,Y_q)$ gives
the upper total mass; integration over $[1,2]$ gives its
lower bound for small $q$. If $\xi>w$ then
$y>(w/(C_hh_A))^{1/b}$. Integrating the upper density to
infinity, and using the total mass for $w\lesssim h_A$, proves
\eqref{an02sepv1:tailupper}.

To specify the interior lower range, let $b_-=\min b>0$ on
the box and choose
$A_*=\max\{2,(2/c_h)^{1/b_-}\}$.
For $w\ge2h_A$ the interval
\[
 A_*(w/h_A)^{1/b}\le y\le2A_*(w/h_A)^{1/b}
\]
has $\xi>w$ throughout. It lies below $Y_q$ for $w\le z_0$,
where one permissible fixed choice is
\begin{equation}
 z_0=\min\left\{1,\frac12
    \min_{\rho\in[13/20,7/10],\ A\in[c_A,C_A]}
                      A^d\left(\frac K{4A_*}\right)^b\right\}>0.
 \label{an02sepv1:lowerrange}
\end{equation}
Indeed the exponent identity
\begin{equation}
 10d-2s+b(\gamma-1)=0
 \label{an02sepv1:cutoffidentity}
\end{equation}
gives $h_A(Kq^{\gamma-1})^b=A_q^dK^b$.
The minimum in \eqref{an02sepv1:lowerrange} is positive by
continuity on a compact parameter set, not by numerical
evaluation. Integrating the lower density on this interval
proves \eqref{an02sepv1:taillower}, with $C_1=2$ (or any
larger fixed constant). The integral factor
$(1-2^{1-a_2})/(a_2-1)$ is uniformly positive.
This argument also displays why the lower estimate must stop
below the endpoint cutoff.

The tail exponent equals $(a_2-1)/b=p_2$.
Q2E's analytic box estimate $P_\rho<153/250$ implies
\[
 \lambda(3-P_\rho)>
 \frac{169}{400}\left(3-\frac{153}{250}\right)
 =\frac{100893}{100000}>1;
\]
therefore $p_2<1$, while $\lambda<2/3$ gives $p_2>0$.
All powers of $A_q$ can be read without absorbing them:
in the power-tail range its coefficient is
\begin{equation}
 B_q h_A^{p_2}
 =D_h^{-2}A_q^{dp_2-1}q^{(10d-2s)p_2}
 =D_h^{-2}A_q^{-3\lambda/2}q^{(10d-2s)p_2}.
 \label{an02sepv1:Apower}
\end{equation}
The profile and radial comparison constants remain two-sided
fixed constants; this identity does not assert an exact tail
prefactor. This completes the theorem.
\end{proof}

Take LCRC's $\Sigma$ to be this entire fixed group, and take
$\mathcal D=\emptyset$, as in roadmap v1.7. The entry cap,
raw-mass budget and upper tail supply exactly
\texttt{an03lcrcv1:profile}, including the joint displacement
and same-field reference required by BULK Lemma 5.1 and
Proposition 5.2. BULK's entry moment lemma can consequently
be applied to this measure. Q2EC and Q2EP supply the separate
primary and whole-Q2 entry contracts; they were not changed.
The reference has entry coordinate bound $R_\Sigma=79/51$.

Nothing here supplies \texttt{an03lcrcv1:distance} or
\texttt{an03lcrcv1:gain} on a later stopped interval, nor
later boundedness of the reference. Those remain S2 obligations
for this same fixed $\Sigma$, at every intermediate endpoint.
In particular BULK's secondary energy corollary remains
conditional on its later distance/gain and clock inputs.
With $\mathcal D=\emptyset$ its separate chart-retention budget
is vacuous; retention of other groups is not established here.
No estimate beyond $t_h$ has been proved or invoked.

\begin{thebibliography}{9}\small
\bibitem{Q2E} Q2E, \path{AN02_Q2_boundary_layers_and_cohort_energy_v1.tex},
\texttt{inc:AN02:q2completion:v1:profile}, crossing and axis-arrival formulas.
\bibitem{PROFILE} PROFILE,
\path{AN02_strong_entry_profile_and_tail_budget_v1.tex},
exact root, uniform arrival, gate integral, Proposition 4.1.
\bibitem{OWPT} OWPT,
\path{AN02_overlap_window_profile_transfer_v1.tex},
\texttt{an02owptv1:jacobian} and \texttt{tangents}; P's global early comparison as used there.
\bibitem{MOMC} MOMC,
\path{AN02_moving_overlap_M_clock_passage_v1.tex},
\texttt{an02momcv1:main}, \texttt{reduction}, \texttt{jacobian},
mass and mean-integral estimates.
\bibitem{Q2EP} Q2EP,
\path{AN02_Q2_entry_through_strong_passage_v1.tex},
whole-Q2 energy theorem.
\bibitem{Q2EC} Q2EC, \path{AN02_Q2_entry_contracts_v1.tex},
\texttt{an02q2ecv1:main}, \texttt{offtransport},
\texttt{delta}, \texttt{targetcontracts}; read from the updated repository.
\bibitem{BULK} BULK,
\path{AN03_bulk_clock_and_secondary_Q2_energy_v1.tex},
Lemma 5.1, Proposition 5.2, secondary corollary.
\bibitem{LCRC} LCRC,
\path{AN03_ledger_cap_residual_closure_v1.tex}, Section 1,
\texttt{an03lcrcv1:profile}, \texttt{distance}, \texttt{gain}.
\bibitem{Roadmap} \path{UNIFIED_ANALYTIC_THEORY_ROADMAP_v1.md},
version 1.7, S0 retained inputs and S1 step 7.
\end{thebibliography}
\end{document}
